\documentclass[a4paper]{article}
% Español (México) con XeLaTeX: fontspec para fuentes Unicode, polyglossia
% para la división silábica y los nombres del documento en la variante mexicana.
\usepackage{fontspec}
\usepackage{polyglossia}
\setdefaultlanguage[variant=mexican]{spanish}
% Los nombres chinos van como «pinyin (original)»: xeCJK da a los caracteres
% Han su propia fuente; sin ella XeLaTeX los omite con un aviso.
% FandolSong no cubre algunos caracteres de nombres (U+73FA); WenQuanYi Zen Hei sí.
\usepackage[AutoFallBack=true]{xeCJK}
\setCJKmainfont{FandolSong-Regular.otf}
\setCJKfallbackfamilyfont{\CJKrmdefault}{WenQuanYi Zen Hei}
% polyglossia no traduce los nombres de listings.
\AtBeginDocument{\providecommand{\lstlistingname}{}\renewcommand{\lstlistingname}{Listado}%
  \providecommand{\lstlistlistingname}{}\renewcommand{\lstlistlistingname}{Índice de listados}}
\usepackage{amsmath, amssymb}
\usepackage{graphicx}
\usepackage[margin=2.35cm]{geometry}
\usepackage[most]{tcolorbox}
\usepackage{booktabs}
\usepackage{tabularx}
\usepackage{longtable}
\usepackage{array}
\usepackage{float}
\usepackage{xcolor}
\usepackage{hyperref}
\usepackage{fancyhdr}
\setCJKmainfont[AutoFakeBold=2.4]{AR PL UMing CN}
\setCJKsansfont[AutoFakeBold=2.4]{Droid Sans Fallback}
\setCJKmonofont[AutoFakeBold=2.4]{Droid Sans Fallback}
\hypersetup{colorlinks=true, linkcolor=blue!55!black, urlcolor=blue!60!black}
\graphicspath{{./}{figures/}}
\setlength{\parskip}{0.55em}
\setlength{\parindent}{2em}
\setlength{\emergencystretch}{3em}
\renewcommand{\arraystretch}{1.18}
\newtcolorbox{knowledgebox}[1]{enhanced, colback=blue!5!white, colframe=blue!70!black, colbacktitle=blue!70!black, coltitle=white, fonttitle=\bfseries, title={#1}, sharp corners, breakable}
\newtcolorbox{importantbox}[1]{enhanced, colback=yellow!10!white, colframe=yellow!75!black, colbacktitle=yellow!75!black, coltitle=black, fonttitle=\bfseries, title={#1}, sharp corners, breakable}
\newtcolorbox{warningbox}[1]{enhanced, colback=red!5!white, colframe=red!70!black, colbacktitle=red!70!black, coltitle=white, fonttitle=\bfseries, title={#1}, sharp corners, breakable}
\pagestyle{fancy}
\fancyhf{}
\fancyfoot[C]{\thepage}
\fancyfoot[R]{\footnotesize\color{gray}\href{https://github.com/hqhq1025/ai-course-notes}{AI Course Notes}}
\renewcommand{\headrulewidth}{0pt}
\renewcommand{\footrulewidth}{0pt}
\newcommand{\videourl}{https://www.youtube.com/watch?v=k82iFzvKFCQ}
\newcommand{\repourl}{https://github.com/hqhq1025/ai-course-notes}

\begin{document}
\begin{titlepage}
\centering
{\Large Notas didácticas de una entrevista a fondo\par}
\vspace{1.0cm}
{\huge\bfseries Apostar por OpenAI, Robinhood y las cuentas del capital en IA\par}
\vspace{0.5cm}
{\Large Zhang Xiaojun conversa con Freda, socia de Altimeter\par}
\vspace{0.35cm}
{\large Elaborado a partir del video público de Zhang Xiaojun Podcast, una transcripción hecha con GPU, la estructura de capítulos y diagramas conceptuales propios\par}
\vspace{0.3cm}
{\large 2025-12-16\par}
\vspace{0.85cm}
\includegraphics[width=0.88\textwidth,height=0.42\textheight,keepaspectratio]{cover.jpg}\par
\vfill
\begin{tcolorbox}[width=0.92\textwidth, colback=black!2!white, colframe=black!55, sharp corners]
\textbf{Título del video}: 125. Conversación con Freda, socia de Altimeter: la apuesta por OpenAI, la historia de Robinhood, los «chicos malos» del capital estadounidense, los cálculos y la burbuja\par
\textbf{Canal del video}: Zhang Xiaojun Podcast\par
\textbf{Duración del video}: 01:24:26\par
\textbf{Enlace del video}: \href{\videourl}{\nolinkurl{\videourl}}\par
\textbf{Nota sobre los materiales}: YouTube no ofrece subtítulos disponibles; el texto se elaboró a partir de una transcripción con faster-whisper large-v3 en CUDA, los capítulos del video, la portada y figuras didácticas propias. Las tomas fijas de la entrevista no se reutilizan en el texto.\par
\textbf{Página del proyecto}: \href{\repourl}{\nolinkurl{\repourl}}\par
\end{tcolorbox}
\end{titlepage}

\tableofcontents
\newpage
\section{Introducción: la IA y el nuevo precio de internet desde la perspectiva del capital}

Esta sección explica primero por qué este episodio debe leerse según «cómo el capital le saca las cuentas a la IA», y no como una lista de inversiones extraída una por una. Con esta organización, el lector puede ubicar a OpenAI, Robinhood, Agentic Commerce y la AI bubble en un mismo mapa de ingresos, costos y control del punto de acceso.

La invitada de este episodio, Freda, es socia de Altimeter Capital, y su perspectiva abarca los mercados primario y secundario. Este episodio se complementa con el reporte trimestral de Guang Mi (广密) en EP127: EP127 se parece más a una revisión de los paradigmas tecnológicos e industriales, mientras que este episodio muestra cómo el capital le «saca las cuentas» a la IA, a las finanzas por internet y a los ejes principales de la bolsa estadounidense para los próximos tres años. El programa trae muchas cifras, pero su valor didáctico no está en que el lector memorice una valuación concreta, sino en entender cómo el mercado de capitales traduce la narrativa tecnológica a ingresos, flujo de efectivo, participación de mercado, capacidad organizacional, personalidad de los fundadores y su expresión en el mercado secundario.

Estas notas no constituyen ninguna recomendación de inversión. El texto conserva las cifras generales y los marcos del programa para ayudar al lector a entender por qué OpenAI se considera una empresa de producto y no solo una empresa de modelos; por qué el flujo de efectivo de las empresas de modelos se parece a una «bola de nieve negativa»; por qué Robinhood, aun con una mano de malas cartas, todavía puede generar Alpha; por qué Agentic Commerce podría reescribir la publicidad, el comercio electrónico y los pagos; y por qué, para juzgar si hay una AI bubble, no basta con preguntar por el ánimo del mercado, sino que hay que observar si los modelos siguen mejorando, si los ingresos se materializan y si el ROIC mejora.

\begin{importantbox}{Tesis central del episodio}
La narrativa del capital sobre la IA está pasando de «quién invierte más» a «quién convierte la inversión en ingresos». El mercado de capitales ya no premia solo el CAPEX y las grandes visiones, sino que pregunta: si las empresas de modelos pueden convertir ChatGPT/API/Agent/publicidad en una pila de ingresos, si las aplicaciones de IA pueden pasar de la bolsa de ingresos electrónicos a la bolsa de gasto en mano de obra, y si las industrias tradicionales pueden reflejar en sus estados financieros la mejora de eficiencia que aporta la IA.
\end{importantbox}

\begin{knowledgebox}{Nota sobre la estrategia visual}
Este video es una entrevista con toma fija, sin slides, pizarrón ni demostraciones de producto. El texto usa la portada para identificar la fuente, y todas las imágenes del cuerpo son diagramas conceptuales sencillos que muestran el marco de inversión, la pila de ingresos, el ciclo virtuoso del negocio y el tablero para evaluar la burbuja; las explicaciones detalladas se encuentran en el texto y en los recuadros de Lectura de la figura.
\end{knowledgebox}

\subsection{Resumen de la sección}

EP125 es una nota de observación del mercado de capitales. Su eje no es «recomendar qué comprar», sino descomponer el modelo de negocio de las empresas de IA y de internet en variables discutibles: de dónde vienen los ingresos, cuándo se desacelera el costo, quién tiene poder de fijación de precios, quién tiene capacidad de ejecución organizacional y quién puede convertir una tendencia en resultados financieros cuantificables.

\section{2020--2025: cómo el ánimo de Silicon Valley pasó de los despidos al predominio de la IA}

Esta sección revisa primero cómo cambió el ánimo del mercado entre 2020 y 2025, porque los juicios de inversión de Freda no surgieron de la nada: se formaron a lo largo de varios ciclos sucesivos de tasas de interés, inflación, despidos, financiamiento de IA y volatilidad en el mercado secundario.

Este capítulo comienza por establecer el contexto temporal. Freda repasa las palabras clave de Silicon Valley entre 2020 y 2025: en 2020, la pandemia y los recortes de tasas; en 2021, un mercado que siguió en auge, aunque las empresas de pequeña y mediana capitalización tocaron techo primero; en 2022, la inflación y la fuerte caída del Nasdaq; en 2023, la IA se convirtió de pronto en el eje del mercado antes de que este alcanzara a reaccionar; 2024 fue el año en que la narrativa de la IA resultó más fácil de sostener, y en 2025 se sumaron los aranceles, la relación entre China y Estados Unidos y una ola de IA que siguió acumulando fuerza.

La observación más importante aquí es que las señales de financiamiento de la IA se adelantaron a la narrativa general. En 2022 el mercado secundario se veía sombrío, pero empresas de modelos como Inflection AI, Anthropic, Cohere y Stability todavía lograban levantar montos inusualmente grandes en sus rondas Serie A. La explicación de Freda es directa: Scaling Law, la compra de GPU y el gasto de capital fueron conceptos que hicieron que el dinero de los VC fluyera muy pronto hacia Nvidia y hacia las empresas de modelos. Aquí aparece la ventaja de los fondos Crossover: si por el momento no se puede saber cuál de las decenas de empresas de modelos saldrá ganadora, se puede expresar una postura sobre la demanda total de IA a través de Nvidia en el mercado secundario.

\begin{knowledgebox}{Términos clave: Crossover, Beta y Alpha}
\begin{tabularx}{\textwidth}{p{0.22\textwidth}p{0.32\textwidth}X}
\toprule
Término & Explicación breve & Papel en este episodio \\
\midrule
Fondo Crossover & Fondo que invierte a la vez en empresas privadas del mercado primario y en empresas cotizadas del mercado secundario & Puede usar el mercado público (Nvidia, Meta, Robinhood, entre otras) para expresar su juicio sobre las tendencias de frontera. \\
Beta & Rendimiento que proviene del alza general del mercado o del sector & Si solo se apuesta por el ciclo de las criptomonedas, comprar Coinbase no necesariamente aporta más Alpha que comprar bitcoin. \\
Alpha & Rendimiento propio de la empresa que va más allá de lo que explican el sector o el índice & Robinhood sirve para mostrar cómo la capacidad de ejecución de una empresa y la revaloración de su negocio generan Alpha. \\
CAPEX & Capital Expenditure, gasto de capital & Los centros de datos de IA, las GPU, el entrenamiento y la infraestructura en la nube dependen en gran medida del CAPEX. \\
\bottomrule
\end{tabularx}
\end{knowledgebox}

\subsection{Resumen de la sección}

Entre 2020 y 2025, Silicon Valley no tuvo una recuperación gradual, sino una serie de cambios bruscos y continuos. Antes de que la IA se volviera el eje principal, el capital ya había detectado señales inusuales en las grandes rondas de financiamiento de las empresas de modelos y en la demanda de Nvidia. Solo con este contexto se entiende por qué este episodio pasa constantemente del mercado primario al secundario y viceversa.
\section{Tres ejes del sector tecnológico en la bolsa estadounidense: AI, reindustrialización y digitalización financiera}

Después de establecer el contexto temporal, esta sección pasa a la visión de conjunto de Freda sobre los ejes actuales del sector tecnológico en la bolsa estadounidense. Ella considera que hoy hay tres ejes muy claros: AI, Re-industrialization (reindustrialización) y Digitization of Finance (digitalización de la industria financiera). Estos tres ejes no son islas paralelas, sino que están conectados entre sí: la AI necesita centros de datos, energía y chips; la reindustrialización dirige el capital hacia la infraestructura dentro de Estados Unidos; y la digitalización financiera converge con la AI en las monedas estables, los pagos y el Agent Commerce.

\begin{figure}[H]
\centering
\includegraphics[width=0.96\textwidth]{us-tech-three-lines.png}
\caption{Los tres ejes del sector tecnológico en la bolsa estadounidense: la AI, la reindustrialización y la digitalización financiera encajan entre sí. Diagrama conceptual propio, elaborado a partir de la conversación de 00:08:12--00:10:20.}
\end{figure}

\begin{knowledgebox}{Lectura de la figura: los tres ejes no son tres mesas de juego independientes}
La AI necesita chips, nube y centros de datos; la reindustrialización aporta inversión en energía, manufactura e infraestructura; la digitalización financiera entra en la cadena de transacciones mediante las monedas estables, los pagos y el Agent Commerce. Las políticas públicas, los aranceles y las elecciones intermedias determinan la velocidad y la dirección de estos flujos de capital.
\end{knowledgebox}

\subsection{Por qué la AI y la reindustrialización están vinculadas}

Esta sección lleva la AI del relato del software de vuelta a la infraestructura física. Antes se dijo que los tres ejes encajan entre sí; el vínculo más directo son los centros de datos y la inversión en energía: el avance de los modelos requiere capacidad de cómputo, y la capacidad de cómputo requiere electricidad, terrenos, chips y capacidad de construcción.

La capa superficial de la AI son los modelos y las aplicaciones; la capa de fondo son la electricidad, los terrenos, los chips, los centros de datos, las redes y la ingeniería de construcción. Freda menciona que los compromisos de inversión en Estados Unidos de países como Japón y Corea del Sur se destinarán a proyectos de infraestructura como energía y centros de datos, y ese dinero, a su vez, sostendrá la construcción de capacidad de cómputo para AI. Es decir, la AI no es un ciclo puramente de software, sino un ciclo de gasto de capital que arrastra consigo la tecnología dura, la energía, la manufactura y los acuerdos diplomáticos.

\begin{warningbox}{Advertencia de interpretación: las empresas de AI no son SaaS tradicional}
El SaaS tradicional suele tener pocos activos, márgenes brutos altos y costos marginales bajos; las empresas de modelos de AI y parte de las aplicaciones de AI requieren, en cambio, inversión continua en inferencia, entrenamiento, datos e infraestructura. Aplicar directamente a las empresas de AI la plantilla de márgenes brutos del software anterior impide ver el cambio en la estructura de costos.
\end{warningbox}

\subsection{Por qué la digitalización financiera vuelve a ser importante}

A continuación se pasa de la electricidad y los chips al dinero y las transacciones. La digitalización financiera vuelve a ser importante porque, si un Agent realiza compras en nombre del usuario, los pagos, las monedas estables, las autorizaciones y los puntos de entrada a las transacciones pasan a formar parte de la comercialización de la AI.

En este episodio, la digitalización financiera no es la fintech en términos genéricos, sino la combinación de monedas estables, pagos, mercados de predicción, Robinhood y Agent Commerce. Freda menciona que los pagos con monedas estables son cheaper/faster/better, y que en el futuro, cuando un Agent pague y compre cosas en nombre del usuario, los conceptos de token o de moneda estable volverán a plantearse. Este eje conecta los puntos de entrada de la AI, el pago y la liquidación, las transacciones en internet y la regulación financiera.

\begin{importantbox}{Implicaciones financieras del Agent Commerce}
Si el Agent se convierte en el punto de entrada de las compras, no solo afectará a la publicidad y al comercio electrónico, sino también a las billeteras de pago, la autorización de transacciones, las monedas estables, el control de riesgos y la captación de clientes de los comercios. El punto de conexión entre la digitalización financiera y la AI es «quién completa la transacción en nombre del usuario».
\end{importantbox}

\subsection{El lugar del mercado chino a los ojos de los inversionistas estadounidenses}

En el programa, la valoración del mercado chino es bastante mesurada: cada vez más inversionistas estadounidenses visitan China y quedan impresionados por los EV y los robots, pero el monto real de inversión no ha cambiado de forma notable. Lo que más les importa son dos cosas: el resultado de las negociaciones entre China y Estados Unidos, y el dinamismo de las IPO en la bolsa de Hong Kong. Dicho de otro modo, al capital estadounidense no le falta interés, sino que está esperando señales sobre las políticas, las salidas y la calidad de las IPO tecnológicas.

\subsection{Resumen de la sección}

Los tres ejes amplían la AI desde las empresas de modelos hacia la energía, la manufactura, los pagos y las políticas públicas. Esto indica al lector que el ciclo de la AI no es un ciclo único de software, sino un ciclo que redistribuye el gasto de capital, la infraestructura industrial y los puntos de entrada financieros.
\section{OpenAI: empresa de producto, bola de nieve negativa y cuatro pilares de ingresos}

La sección anterior expuso el hilo macroeconómico; esta entra en el caso central, OpenAI. El juicio decisivo de Freda cuando apostó por OpenAI fue haber visto antes que nadie que se trataba de una empresa de producto y no solo de una empresa de modelos. En 2023 el mercado todavía temía que ChatGPT no lograra retener a sus usuarios, y el panorama competitivo con Perplexity, Claude, Grok y otros tampoco era claro. Aquí Perplexity es el nombre de una empresa de búsqueda y respuesta a preguntas con IA, no la perplexity/PPL de la evaluación de modelos de lenguaje; esta última es la métrica de «perplejidad», que se obtiene al aplicar la exponencial a la entropía cruzada. La ventaja de haber llegado primero, la experiencia de usuario y el lugar que ChatGPT ocupa en la mente de los usuarios le dieron a OpenAI un valor como punto de acceso distinto del de un laboratorio de modelos común.

\subsection{Por qué se dice que OpenAI es una empresa de producto}

Esta sección explica primero la tesis contraria al consenso de «empresa de producto», porque determina que el ancla de la valuación de OpenAI no sea solo la capacidad del modelo, sino también el punto de acceso de los usuarios, la retención, los flujos de trabajo empresariales y las futuras interfaces de monetización mediante publicidad y comercio electrónico.

Si una empresa de modelos es solo una API o un laboratorio, su valor proviene de la capacidad; una empresa de producto, en cambio, posee además la relación con el usuario, la marca, la retención, el punto de acceso y la interfaz de monetización. ChatGPT no es un producto con fuertes efectos de red como los de Meta; se parece más a una distribución uno a uno. Sin embargo, si los usuarios lo adoptan como asistente personal, asistente de trabajo y punto de acceso empresarial, puede generar una retención muy fuerte. Freda subraya en particular que el mercado subestima el segmento empresarial: la proporción entre ingresos de usuarios individuales y de empresas no está tan cargada hacia el consumidor como se imagina desde fuera, y cuando los usuarios empresariales conectan el correo, los sistemas internos y las suites de oficina, la idea de un gran punto de acceso único se vuelve todavía más clara.

\begin{importantbox}{Diferencia entre empresa de producto y empresa de modelos}
Una empresa de modelos responde a «¿la capacidad es suficientemente fuerte?»; una empresa de producto responde a «¿el usuario está dispuesto a seguir usándolo, a pagar, a entregarle su contexto y a dejar que entre en su flujo de trabajo?». La tesis central de OpenAI, contraria al consenso, es que ChatGPT convirtió la capacidad del modelo en un punto de acceso de producto que puede distribuirse, retener usuarios y comercializarse.
\end{importantbox}

\subsection{Bola de nieve negativa: la aritmética del flujo de efectivo de una empresa de modelos}

Esta sección explica el pasaje de aritmética del modelo de negocio con mayor valor didáctico del programa. Freda describe a la empresa de modelos como una «bola de nieve negativa»: si el costo de entrenamiento de la generación anterior es 1, al año siguiente los ingresos pueden recuperar 2, pero el costo de entrenamiento de la siguiente generación puede llegar a 10, de modo que el flujo de efectivo es \(+2-10=-8\). Mientras el costo de entrenamiento siga multiplicándose, el consumo de efectivo será cada vez mayor.

\begin{figure}[H]
\centering
\includegraphics[width=0.96\textwidth]{openai-negative-flywheel.png}
\caption{Bola de nieve negativa de OpenAI: los ingresos persiguen el entrenamiento de la generación anterior, mientras el entrenamiento de la siguiente generación sigue ampliando el gasto. Diagrama conceptual de elaboración propia, basado en la conversación de 00:12:14--00:16:45.}
\end{figure}

\begin{knowledgebox}{Lectura de la figura: el punto de inflexión surge de un menor crecimiento del entrenamiento o de un mayor múltiplo de ingresos}
La figura no afirma que OpenAI nunca podrá ganar dinero; muestra las condiciones para que el flujo de efectivo se vuelva positivo: o bien los ingresos del segundo año dejan de ser solo el doble del costo de la generación anterior y son mayores, o bien el costo de entrenamiento del siguiente modelo deja de crecer en órdenes de diez veces. El verdadero punto de inflexión llega cuando se desacelera el crecimiento de la inversión en modelos, cuando mejora el ROI, o cuando ocurren ambas cosas a la vez.
\end{knowledgebox}

Freda compara OpenAI con Netflix: en sus primeros años, Netflix también tuvo flujo de efectivo negativo durante mucho tiempo, porque debía invertir continuamente en contenido; cuando el crecimiento de esa inversión disminuyó o ya no pudo seguir produciendo, el flujo de efectivo se volvió positivo de golpe. Con las empresas de modelos ocurre algo similar: la etapa más difícil es aquella en que crecen la capacidad y los usuarios, pero también el costo de entrenamiento; si algún día ese costo deja de crecer de forma exponencial, el estado de resultados se verá muy bien. Esto no significa dejar de entrenar por completo, sino que el crecimiento pasa de ser explosivo a ser de mantenimiento.

\begin{warningbox}{Límites de la analogía con Netflix}
El contenido de Netflix no es del todo exclusivo y los usuarios pueden suscribirse a varias plataformas al mismo tiempo; ChatGPT se parece más a un punto de acceso de búsqueda y de asistente personal, por lo que la concentración del mercado podría ser mayor. Aun así, ambos casos muestran que entre la inversión capitalizada inicial y la mejora posterior del flujo de efectivo puede haber un desfase muy largo.
\end{warningbox}

\subsection{Los cuatro pilares de ingresos de OpenAI}

Ya se explicó por qué el flujo de efectivo se ve mal a corto plazo; esta sección pasa al lado de los ingresos: si OpenAI quiere sostener una inversión enorme, no puede depender de un solo producto de suscripción individual, sino que debe apilar varios pilares de ingresos.

En el programa, la estructura de ingresos de OpenAI se divide en cuatro líneas: ChatGPT, API, Agent y nuevos productos/publicidad. Actualmente ChatGPT es la mayor parte e incluye suscripciones individuales y empresariales; la API se orienta al consumo por parte de desarrolladores y empresas; Agent incluye alianzas como la de SoftBank, la ejecución de tareas y los flujos de trabajo futuros; la publicidad y el comercio electrónico corresponden a la monetización de los usuarios gratuitos y a la idea de un punto de acceso para búsqueda y compras.

\begin{figure}[H]
\centering
\includegraphics[width=0.96\textwidth]{openai-four-pillars.png}
\caption{Los cuatro pilares de ingresos de OpenAI: ChatGPT, API, Agent, y publicidad/nuevos productos junto con el punto de acceso empresarial. Diagrama conceptual de elaboración propia, basado en la conversación de 00:16:45--00:20:49.}
\end{figure}

\begin{knowledgebox}{Lectura de la figura: no mirar solo las suscripciones individuales de ChatGPT}
En la figura, ChatGPT es hoy la mayor parte, pero la API, Agent, el punto de acceso empresarial y la publicidad/comercio electrónico determinarán el techo de ingresos. El segmento empresarial es especialmente importante porque el modelo de cobro del software es más directo: a las suites de oficina de terceros no necesariamente les importa estar en la parte más alta del embudo de adquisición de clientes; mientras los usuarios sigan pagando, el punto de acceso de Agent puede incluso aumentar la profundidad de uso del software.
\end{knowledgebox}

\subsection{Competencia, IPO y rendimiento de la inversión}

La competencia de OpenAI proviene de Google, Anthropic, Meta, xAI, los modelos de código abierto y los futuros Neo Labs. La postura de Freda es que la ventaja de OpenAI es muy sólida, pero el rendimiento de la inversión no puede medirse solo por el aumento de la valuación de 30,000 millones a 500,000 millones, porque en el mercado primario la dilución accionaria es muy fuerte y lo que el inversionista realmente obtiene es el crecimiento del valor por acción. También considera que, si OpenAI saliera a bolsa alrededor de 2027 con una capitalización de 1 billón de dólares, eso no equivaldría necesariamente a la cima de una burbuja, siempre que sus ingresos alcancen una escala acorde; Netscape salió a bolsa en 1995 y la burbuja de internet no estalló realmente sino varios años después.

\begin{warningbox}{El múltiplo de valuación no equivale al rendimiento real del inversionista}
Los medios suelen decir cuántas veces se multiplicó la valuación de una empresa, pero en el mercado primario las acciones se diluyen con cada ronda de financiamiento. El inversionista debe fijarse en el precio por acción, las condiciones, la preferencia de liquidación y las rondas posteriores, no solo en la valuación total.
\end{warningbox}

\subsection{Resumen de la sección}

El relato financiero de OpenAI se compone de tres elementos: es una empresa de producto que posee el punto de acceso de ChatGPT; su flujo de efectivo se comporta a corto plazo como una bola de nieve negativa, cuyo punto de inflexión depende del costo de entrenamiento y del ROI de los ingresos; y su conjunto de ingresos no se limita a las suscripciones individuales, sino que también incluye la API, Agent, el punto de acceso empresarial y la publicidad/comercio electrónico.
\section{Anthropic, Neo Labs y la diferenciación entre empresas de modelos}

La sección anterior desglosó OpenAI en productos y flujo de efectivo; esta sección examina la diferenciación entre empresas de modelos. La diferencia entre Anthropic y OpenAI no se reduce a que sus modelos tengan nombres distintos: difieren la estructura de clientes, la línea de productos, el control del margen bruto y los supuestos de sus proyecciones financieras. Neo Labs representa otro tipo de diferenciación: equipos de investigación que salen de las grandes empresas de modelos y siguen apostando por Continuous Learning, el entrenamiento personalizado y el ecosistema de código abierto.

\begin{figure}[H]
\centering
\includegraphics[width=0.94\textwidth]{openai-vs-anthropic.png}
\caption{OpenAI vs Anthropic: una empresa de productos 2C frente a una empresa 2B/Coding. Diagrama conceptual propio, elaborado a partir de la conversación de 00:28:25--00:31:25.}
\end{figure}

\begin{knowledgebox}{Lectura de la figura: las dos empresas no son copias del mismo negocio}
La ventaja de OpenAI se inclina hacia el posicionamiento entre los consumidores, el asistente personal y el punto de entrada empresarial; más de la mitad de los ingresos de Anthropic proviene del segmento 2B, y la empresa pone más énfasis en Coding, Finance, los modelos de Agent y los trabajadores del conocimiento en las empresas. Ambas pueden ir muy bien, pero difieren en la calidad de sus ingresos, sus vías de venta y la forma en que organizan sus productos.
\end{knowledgebox}

\subsection{El crecimiento de Anthropic y el control del margen bruto}

Freda menciona que los ingresos anualizados de Anthropic crecieron muy rápido, de 1,000 millones a 7,000 millones, y que en los últimos meses su incremento mensual de ingresos anualizados se ha acercado en magnitud al de OpenAI. Más importante es la mejora del margen bruto: mediante la arquitectura del modelo, el auto scaling, la optimización de la inferencia y otros medios, el margen bruto de una empresa de modelos no está condenado a ser bajo. Ella considera que no es inimaginable que, a largo plazo, el margen bruto de las empresas de modelos llegue a 70\%--80\%.

Aquí hay que atender a la diferencia en los supuestos de las proyecciones financieras. Según ella, la proyección de Anthropic parece suponer que después de 2028 el cómputo total ya no crecerá de forma considerable, de modo que la desaceleración del crecimiento de los costos de entrenamiento eleva el margen de utilidad; la proyección de OpenAI, en cambio, parece suponer que el ROI de los modelos aumenta año con año, es decir, que cada peso invertido en entrenamiento podrá recuperarse en un múltiplo mayor en el futuro. Ninguno de los dos supuestos es un hecho en sí mismo: son modelos sobre los que conviene que los inversionistas pregunten a fondo.

\begin{importantbox}{Las dos palancas de las proyecciones financieras de una empresa de modelos}
La primera es la palanca de costos: si el crecimiento de los costos de entrenamiento e inferencia se desacelera. La segunda es la palanca de ingresos: si cada generación de modelos aporta un ROI mayor. Cada empresa puede construir su relato de mejora del margen de utilidad con una palanca distinta.
\end{importantbox}

\subsection{Neo Labs: la lógica de financiamiento de la siguiente generación de empresas de modelos}

Esta subsección amplía la perspectiva de los dos líderes, OpenAI y Anthropic, a los nuevos laboratorios. Si, como se dijo antes, las empresas de modelos se diferenciarán, la pregunta de Neo Labs es: cuando las empresas de modelos punteras ya son muy fuertes, ¿en qué pueden apoyarse los nuevos equipos para obtener financiamiento y diferenciarse?

En el programa se mencionan nombres conocidos como SSI y Thinking Machines, y también que en los últimos meses varias Neo Labs más han obtenido financiamiento; muchas fueron fundadas por personas con trayectoria académica y los montos recaudados son muy grandes. La duda de Freda es: ¿los modelos grandes terminarán por concentrarse en una o dos empresas punteras, mientras las demás se vuelcan al código abierto? El primer producto de Thinking Machines, Thinker, se describe como una herramienta que ayuda a los usuarios a hacer entrenamiento personalizado de modelos de código abierto; en el fondo, es una apuesta a que el mercado de código abierto será cada vez más grande.

\begin{knowledgebox}{Términos clave: Neo Labs y personalización de código abierto}
\begin{tabularx}{\textwidth}{p{0.24\textwidth}p{0.32\textwidth}X}
\toprule
Término & Explicación breve & Relación con este episodio \\
\midrule
Neo Labs & Nueva generación de laboratorios de modelos o de investigación, a menudo fundados por talento proveniente de grandes empresas o de la academia & Aprovechan las oportunidades del siguiente paradigma, la personalización de código abierto y los modelos especializados. \\
Continuous Learning & Aprendizaje continuo: mejorar el modelo a partir de nuevas interacciones y datos & En el programa se menciona que SSI podría estar explorando esta dirección. \\
Entrenamiento personalizado de código abierto & Entrenamiento específico para el usuario sobre la base de un modelo de código abierto & Thinker, de Thinking Machines, se usa para ilustrar esta dirección. \\
Valuación con OpenAI como referencia & Usar la valuación total de OpenAI como ancla relativa para otras empresas de modelos & Referencia psicológica habitual en el financiamiento de empresas de modelos en el mercado primario. \\
\bottomrule
\end{tabularx}
\end{knowledgebox}

\subsection{Resumen de la sección}

Anthropic y Neo Labs muestran que las empresas de modelos ya se han diferenciado: unas son el punto de entrada para los consumidores, otras se enfocan en empresas y Coding, otras en la personalización de código abierto y otras apuestan por el siguiente paradigma. El mercado de capitales seguirá usando a OpenAI como ancla, pero las diferencias reales provienen de la estructura de clientes, el control del margen bruto y la línea de productos.
\section{Robinhood: cómo obtener Alpha a partir de una mala mano}

Después de varias secciones dedicadas a empresas de modelos, esta sección pasa a Robinhood para mostrar cómo el mismo método de análisis de capital se aplica a una empresa de finanzas por internet ajena a la IA. Aquí lo importante no son las subidas y bajadas de la acción, sino cómo separar, dentro de un negocio cíclico, las variables que la empresa sí puede controlar.

Las secciones anteriores giraron en torno a la IA; esta pasa a Robinhood. El caso es importante porque muestra cómo analiza Freda una empresa de finanzas por internet: no se fija solo en la Beta del sector, sino en qué variables puede controlar la empresa. Según ella, el negocio de base de Robinhood no es bueno, porque la intermediación bursátil es muy cíclica y los ingresos por operaciones fluctúan con la actividad del mercado. Sin embargo, mediante la diversificación, la ganancia de participación de mercado, el poder de fijación de precios, el control de costos y la estructura organizacional, la empresa intenta obtener Alpha a partir de una mala mano.

\begin{figure}[H]
\centering
\includegraphics[width=0.96\textwidth]{robinhood-alpha-loop.png}
\caption{Trayectoria de Alpha de Robinhood: del negocio cíclico a la diversificación, la participación de mercado, el poder de fijación de precios y el control de costos. Diagrama conceptual propio, elaborado a partir de la conversación de 00:32:31--00:40:29.}
\end{figure}

\begin{knowledgebox}{Lectura de la figura: Robinhood no sigue la simple lógica de una casa de bolsa en un mercado alcista}
El primer paso de la figura reconoce que el negocio de operaciones es muy cíclico; los cuatro pasos siguientes son variables que la empresa puede controlar: diversificar las líneas de ingresos, ganar participación en las cuentas de usuarios jóvenes, aumentar el poder de fijación de precios en negocios como las criptomonedas y controlar estrictamente los costos operativos. El Alpha proviene de estas variables, no solo de la subida del mercado.
\end{knowledgebox}

\subsection{Beta y Alpha: por qué no basta con comprar el ciclo}

Esta subsección aclara primero una habilidad básica del análisis de inversiones: una misma subida puede deberse al nivel general del sector o a la ejecución propia de la empresa. La primera es Beta; solo la segunda se acerca más al Alpha.

Freda menciona que, si se concluye que llegó un mercado alcista de criptomonedas, comprar Coinbase no necesariamente genera más Alpha que comprar bitcoin; en cambio, Robinhood siguió una trayectoria propia más fuerte que la de bitcoin y la del Nasdaq. El valor didáctico de este juicio es que el inversionista debe saber con honestidad de dónde proviene lo que gana. Si el rendimiento proviene por completo del ciclo del sector, es Beta; solo si proviene de que la empresa mejora continuamente su negocio, gana participación de mercado y modifica su poder de fijación de precios, se acerca más al Alpha.

\begin{warningbox}{No confundir una «subida correlacionada» con capacidad de la empresa}
Que la acción de una empresa suba junto con su sector no significa que la empresa haya generado Alpha. Para juzgar la capacidad de la empresa, hay que comparar con el índice del sector, los precios de los activos clave y la Beta del mercado, y después ver si lo restante proviene de la estrategia y la ejecución de la empresa.
\end{warningbox}

\subsection{Edad de los usuarios de Robinhood y migración de cuentas}

A continuación se pasa de la atribución de rendimientos a la estructura de usuarios, porque el potencial de largo plazo de Robinhood no proviene solo de la frecuencia de operaciones, sino también de que los activos de los usuarios jóvenes migran de forma natural, a medida que envejecen, hacia un mismo sistema de cuentas.

La edad promedio de los usuarios de Robinhood ronda los treinta y tantos años, y a partir de los 35 años la acumulación de riqueza de los estadounidenses se acelera notablemente. La lógica de Freda es que cada generación tiene su propia cuenta: los mayores de 60 años usan Charles Schwab, las personas de mediana edad usan E*Trade, y los jóvenes y quienes están por entrar en la mediana edad probablemente usen Robinhood. Una vez que los activos entran en la cuenta, Robinhood puede aumentar los activos y los ingresos por usuario mediante operaciones, gestión patrimonial, servicios bancarios, mercados de predicción, negocios internacionales y, en el futuro, fondos de VC, entre otras vías.

Por eso ella valora la trayectoria de Robinhood hacia una aplicación financiera integral. La cuenta de operaciones es solo la puerta de entrada; el verdadero valor de largo plazo está en la acumulación de activos, la gestión patrimonial y el supermercado financiero.

\subsection{Estructura organizacional: Single GM y producción rápida de productos}

Esta subsección examina además la organización, porque la diversificación no ocurre solo por escribirla en una presentación de PPT sobre estrategia: se necesita una estructura organizacional que permita a varias unidades de negocio experimentar y corregir errores rápidamente al mismo tiempo.

Freda menciona también que en 2022 Robinhood prácticamente rehízo la empresa y adoptó una estructura Single GM. Cada unidad de negocio cuenta con su propia estructura de PM, Developer, CFO, HR, etc., con el objetivo de dar a las unidades de negocio la capacidad de innovar con rapidez. Este ajuste organizacional explica por qué la cantidad de productos que Robinhood lanza en un año parece equivaler a lo que otras empresas lanzan en cinco a diez años.

\begin{importantbox}{Lo que implica Single GM para el producto}
Cuando una empresa quiere pasar de un único negocio de operaciones a una aplicación financiera con varios negocios, una organización centralizada frena la experimentación. Single GM hace que cada línea de negocio funcione más como una pequeña empresa, capaz de explorar al mismo tiempo oportunidades como servicios bancarios, mercados de predicción, gestión patrimonial, internacionalización y tokenización.
\end{importantbox}

\subsection{Resumen de la sección}

El valor del caso Robinhood está en entrenar la capacidad de «descomponer un negocio». Un negocio muy cíclico y de base imperfecta aún puede obtener Alpha mediante la estructura de usuarios, la velocidad de producto, la diversificación, el control de costos y el diseño organizacional.
\section{Niños rebeldes, erizos y nuevas especies: el taste de los fundadores para el capital de Silicon Valley}

La sección anterior trató el negocio de Robinhood; esta sección trata de los fundadores y del taste del capital. En el programa aparecen varias expresiones muy gráficas: el niño bien portado, el niño rebelde, el tipo erizo, el inconforme y el tipo Nezha (哪吒). No sirven para poner etiquetas, sino para describir cómo el capital de Silicon Valley identifica a las empresas founder-led capaces de «romper con lo establecido».

\begin{figure}[H]
\centering
\includegraphics[width=0.96\textwidth]{founder-archetypes.png}
\caption{Perfiles de fundadores en Silicon Valley: el niño bien portado, el niño rebelde, el erizo, el inconforme, Nezha y Founder-led. Diagrama conceptual de elaboración propia, a partir de la conversación de 00:40:29--00:44:26 y 00:57:13--00:58:22.}
\end{figure}

\begin{knowledgebox}{Lectura de la figura: no es una prueba de personalidad, sino el taste de un inversionista}
El niño bien portado representa el apego a las normas y la prudencia; el niño rebelde, la disposición a romper las reglas; el tipo erizo, una profundidad extrema en un solo punto; el inconforme, el desafío al orden establecido; y el tipo Nezha, no jugar según las reglas del juego. Lo que el capital busca en realidad es la capacidad de ejecución sostenida de una empresa founder-led, no la rebeldía por la rebeldía misma.
\end{knowledgebox}

\subsection{Rasgos comunes de las empresas disruptivas}

Esta subsección pasa de las etiquetas de los fundadores a los fundamentos de la empresa. Que un fundador sea «rebelde» o «inconforme» no tiene valor por sí mismo; ese carácter solo se convierte en un rasgo por el que el capital está dispuesto a apostar cuando se pone al servicio de un mercado grande, un producto sólido y una ejecución rápida.

Freda resume varios rasgos comunes de las empresas disruptivas. Primero, un TAM suficientemente grande; por ejemplo, Robinhood quiere ser la superaplicación de toda la industria financiera. Segundo, un producto suficientemente bueno, idealmente capaz de crecer a gran velocidad sin necesidad de un gran esfuerzo de ventas. Tercero, coincidir en el momento, el lugar y las personas adecuados: la madurez de los teléfonos inteligentes y del ancho de banda dio origen al video corto; el fin de la ley de Moore y las GPU impulsaron la IA; el vacío en los mercados de ciudades pequeñas y zonas rurales dio origen a Pinduoduo (拼多多). Muchos disruptores parecían ilógicos cuando aparecieron: Robinhood con comisiones de cero, OpenAI gastando dinero sin freno, Netflix transformando la renta de DVD, TikTok reescribiendo la distribución de contenido.

\begin{importantbox}{Lo disruptivo no es lo «extraño» en sí}
Que una empresa parezca ilógica es solo una condición necesaria, no suficiente. Además, debe enfrentarse a un mercado enorme, contar con un producto muy sólido, coincidir con el momento tecnológico y social adecuado, y ser capaz, mediante la capacidad de ejecución de su organización, de convertir lo contraintuitivo en el nuevo sentido común.
\end{importantbox}

\subsection{Mercados de predicción: una nueva especie}

A continuación se toma el mercado de predicción como muestra de una «nueva especie». Es a la vez un producto financiero y un mecanismo de descubrimiento de información, y actúa en conjunto con la participación de pequeños inversionistas, una regulación laxa y la crisis de confianza en los medios.

El prediction market (mercado de predicción) que se menciona en el programa es un ejemplo de «nueva especie». Plataformas como Kalshi y Polymarket permiten a los usuarios apostar pequeñas cantidades de dinero sobre eventos: resultados electorales, informes financieros que superan las expectativas, fechas de lanzamiento de modelos e incluso eventos del mundo del entretenimiento. Freda considera que no se trata de un mercado de nicho, porque el volumen de operaciones crece con rapidez y tanto Bloomberg como Google Finance ya empezaron a integrar sus datos.

El contexto en el que surgen los mercados de predicción incluye la crisis de confianza en los principales medios de Estados Unidos, la incertidumbre derivada de la pandemia, los aranceles y la inestabilidad política, una regulación relativamente laxa y la creencia de que «siempre hay alguien que sabe algo». Su relación con Robinhood consiste en que el mercado de predicción es a la vez un producto financiero y una nueva vía para que el público participe en la fijación del precio de los eventos.

\subsection{Resumen de la sección}

Al capital de Silicon Valley no le atraen las «empresas obedientes», sino las que, en un mercado grande, logran formar una nueva especie gracias a un producto sólido, una ejecución firme y un criterio contraintuitivo. El perfil del fundador es solo la envoltura; lo que de verdad importa es el mercado, el producto, el momento y la capacidad de la organización.
\section{Conducción autónoma, robótica y cómo expresar una tesis en el mercado privado y en el mercado público}

Las secciones anteriores trataron las empresas de modelos y la tecnología financiera; esta sección agrega la dirección del mundo físico: la conducción autónoma y la robótica. Freda valora mucho la capacidad de ejecución de Waymo: lanza el servicio en nuevas ciudades cada vez más rápido, amplía sus zonas de operación y sus ingresos empiezan a poder estimarse. Al mismo tiempo, advierte que las empresas de robótica alcanzan con facilidad valuaciones muy altas en el mercado privado (primario), mientras que en el mercado público (secundario) puede ser más sencillo expresar la tendencia mediante Nvidia, Tesla, la cadena de suministro relacionada u otros valores cotizados.

\subsection{Waymo y la comercialización de la conducción autónoma}

En el episodio se menciona que Waymo opera en varias zonas y que el tiempo para lanzar el servicio en una ciudad pasó de varios años en San Francisco a plazos más cortos en Austin y Silicon Valley; además, el costo por ciudad no es tan alto como se pensaba. Lo relevante de este dato es que, una vez que la conducción autónoma pasa de la «demostración tecnológica» a la «expansión operativa», los inversionistas empiezan a calcular el costo de replicar el servicio en cada ciudad, el número de vehículos, los ingresos anualizados, la velocidad de aprobación regulatoria y la economía unitaria.

\begin{knowledgebox}{Variables que hay que observar en la comercialización de la conducción autónoma}
\begin{tabularx}{\textwidth}{p{0.23\textwidth}p{0.33\textwidth}X}
\toprule
Variable & Pregunta que hay que responder & Por qué importa \\
\midrule
Velocidad de apertura de ciudades & Cuánto tarda una nueva ciudad en pasar de pruebas a operación & Determina la velocidad con que se replican los ingresos. \\
Costo por ciudad & Cuánto capital e inversión operativa requiere entrar en una ciudad & Determina si la expansión es sostenible. \\
Tamaño de la flota & Cuántos vehículos prestan servicio estable en las calles & Determina la oferta y la experiencia del usuario. \\
Permisos regulatorios & Si puede circular por autopistas, operar entre zonas y ampliar el área de servicio & Determina los límites del negocio. \\
Economía unitaria & Cuáles son los ingresos, el mantenimiento, el seguro y la depreciación por vehículo y por día & Determina el modelo de rentabilidad a largo plazo. \\
\bottomrule
\end{tabularx}
\end{knowledgebox}

\subsection{Robótica y su expresión en el mercado público}

Esta sección vuelve a situar la conducción autónoma y la robótica en la perspectiva de un fondo Crossover. Antes se dijo que la incertidumbre en el mercado privado es alta; por eso los inversionistas suelen preguntarse: ¿existe una forma más líquida y más transversal de expresar la misma tesis sobre la industria en el mercado público?

La robótica y la inteligencia corporeizada siguen siendo una dirección de largo plazo, pero en el mercado privado la incertidumbre de cada empresa individual es muy alta. Desde la perspectiva Crossover de Freda, la pregunta es: si crees en la robótica o en la infraestructura de IA, no necesariamente tienes que apostar por una sola empresa privada; también puedes expresar esa tesis mediante Nvidia, los chips, la energía, la nube, la cadena de suministro o plataformas que cotizan en bolsa. Este razonamiento se parece al de comprar Nvidia en 2023 para expresar la demanda agregada de las empresas de modelos.

\begin{warningbox}{El mercado privado y el mercado público expresan riesgos distintos}
En el mercado privado se puede obtener el rendimiento explosivo de una sola empresa, pero el capital queda inmovilizado por mucho tiempo, la valuación es poco transparente y la tasa de fracaso es alta. El mercado público tiene más liquidez y permite expresar la línea principal de una industria, pero también es más sensible a los factores macroeconómicos y al sentimiento del mercado. No se trata de que uno sea mejor que otro, sino de que tienen estructuras de riesgo y de horizonte temporal diferentes.
\end{warningbox}

\subsection{El mercado privado depende del consenso; el mercado público, de lo que no es consenso}

A continuación se resumen las diferencias entre el mercado privado y el mercado público. Los casos de Robinhood y Nvidia mostraron que el momento más interesante del mercado público suele llegar cuando el cambio real en la empresa ya ocurrió, pero el mercado todavía no lo ha incorporado por completo en el precio.

En el episodio hay una frase muy concisa: «el mercado privado depende del consenso; el mercado público, de lo que no es consenso». En el mercado privado, las empresas más solicitadas suelen necesitar consenso para levantar grandes rondas de financiamiento; en el mercado público, si todos están de acuerdo, el precio normalmente ya refleja ese consenso. El rendimiento excedente real proviene de entender el cambio en un negocio antes que el mercado, por ejemplo, la capacidad de ejecución de Robinhood y su expansión hacia varias líneas de negocio.

\subsection{Resumen de la sección}

La enseñanza común de la sección sobre conducción autónoma y robótica es que, en la tecnología de frontera, no basta con ver la demo: también hay que analizar las variables de comercialización y la forma de expresar la tesis. Un inversionista puede apostar por una empresa, o bien por la cadena de suministro, los chips, la nube, la energía o valores cotizados más líquidos del mercado público.
\section{Balance de los sectores de IA en 2025: Coding, video, Agentic Commerce y aplicaciones verticales}

Esta sección pasa de los casos de empresas a los sectores. Freda plantea un secreto a voces del mundo de la inversión: si un campo genera 1,000 millones de dólares de ingresos anualizados en un año, hay que estudiarlo con seriedad, se haya entendido antes o no, porque el mercado está indicando que ahí pueden surgir empresas de talla mundial. El caso más representativo de 2025 es Coding, que en un año generó más de 5,000 millones de dólares de ingresos anualizados.

\begin{figure}[H]
\centering
\includegraphics[width=0.96\textwidth]{ai-sector-map-2025.png}
\caption{Mapa de los sectores de IA en 2025: Coding, video, Agentic Commerce, finanzas/derecho, salud/atención al cliente y valuación de aplicaciones de IA. Diagrama conceptual propio, elaborado a partir de la conversación de 01:03:43--01:09:51.}
\end{figure}

\begin{knowledgebox}{Lectura de la figura: el interés por un sector se mide por la velocidad de sus ingresos}
Coding es el sector con el crecimiento explosivo de ingresos más claro; el video puede transformar la industria de los medios porque permite optimizar directamente el contenido; Agentic Commerce afectará la publicidad, el comercio electrónico y los pagos; en las aplicaciones verticales como finanzas, derecho, salud y atención al cliente, lo determinante son el ROI y los datos de cada industria.
\end{knowledgebox}

\subsection{Por qué Coding es el campo más adecuado para que las empresas de modelos lo desarrollen ellas mismas}

Esta subsección analiza primero Coding, porque es la muestra más clara de ingresos de aplicaciones de IA en 2025 y la que mejor muestra qué escenarios verticales absorberán directamente las empresas de modelos.

A diferencia de las finanzas, el derecho o la salud, Coding no requiere recopilar grandes cantidades adicionales de conocimiento del sector: las empresas de modelos ya cuentan, por su propia naturaleza, con datos de código, acceso a los desarrolladores y cadenas de herramientas. Tras el lanzamiento de productos como Claude Code y OpenAI Codex, los ingresos subieron en línea recta, lo que muestra que las empresas de modelos tienen una ventaja considerable en este campo vertical. Sin embargo, el panorama competitivo sigue sin estar claro: si Google u otra empresa de modelos ofreciera gratis sus capacidades de Coding, eso presionaría a empresas independientes como Cursor.

\begin{warningbox}{No todas las aplicaciones verticales son adecuadas para las empresas emergentes}
En la generación anterior, el software solía pasar primero por una capa horizontal general y después a los nichos verticales; en la era de la IA, los grandes modelos absorben la capa general, y muchas oportunidades para empresas emergentes se desplazan hacia escenarios verticales. Pero Coding es la excepción: está demasiado cerca de los datos de entrenamiento y del ecosistema de desarrolladores de las empresas de modelos, así que las empresas independientes deben encontrar una experiencia de producto superior o una barrera de entrada basada en el flujo de trabajo.
\end{warningbox}

\subsection{Video: de la función de emparejamiento al objeto optimizable}

Freda confía en el video porque la industria de los medios es enorme, el video es una forma de entrada de alto ancho de banda y, por primera vez, el video puede ser optimizado directamente por un modelo. En la era de TikTok lo que se optimizaba era el emparejamiento: los creadores producían videos y luego estos se emparejaban con los usuarios; en la era del video generativo, el propio video puede optimizarse según una función objetivo, por ejemplo la participación, la tasa de clics en anuncios, la tasa de conversión o el efecto de marca. Este cambio convertirá el video de un problema de distribución de contenido en un problema de generación y optimización de contenido.

\subsection{Agentic Commerce: la reorganización de la publicidad, el comercio electrónico y los pagos}

Agentic Commerce se refiere a que un Agent compra productos en nombre del usuario: entiende la necesidad, compara productos, recomienda, hace el pedido y paga. Freda considera que para los comerciantes esto podría ser una pesadilla, porque quieren que los usuarios lleguen directamente a su propio sitio web para hacer cross-sell, vender publicidad y controlar el tráfico. Booking.com depende del tráfico que le envía Google y Amazon tiene enormes ingresos publicitarios; ambos modelos solo funcionan si el usuario visita directamente la plataforma.

\begin{importantbox}{Quiénes ganan y quiénes pierden con Agentic Commerce}
Los beneficiados podrían ser los comerciantes de nicho que antes no podían pagar publicidad, porque un Agent puede emparejar con precisión productos de cola larga; los perjudicados podrían ser las plataformas que dependen del tráfico de entrada, de los espacios publicitarios y de los eslabones intermedios de pago. Las empresas de pagos también se revaluarán, porque el propio Agent puede convertirse en la cartera del usuario y en su punto de autorización.
\end{importantbox}

\subsection{Valuación de empresas de aplicaciones de IA: margen bruto o utilidad absoluta}

Esta subsección vuelve al criterio de valuación, porque las aplicaciones de IA no tienen, como el software tradicional, una plantilla natural de márgenes brutos altos y estables. Los inversionistas deben volver a decidir: ¿hay que fijarse en el margen bruto porcentual o en el margen bruto por usuario y el monto absoluto de la utilidad?

El software tradicional es como la «pechuga de pollo»: el margen bruto y la tasa de retención neta de ingresos son muy homogéneos, y los inversionistas pueden valuarlo en serie. Las aplicaciones de IA son distintas: cuanto más se usan, mayor es el costo de inferencia, y el margen bruto puede incluso disminuir. La recomendación de Freda es fijarse en el monto absoluto de la utilidad, porque el monto de los contratos de IA y el margen bruto por usuario pueden ser mucho mayores que en el software tradicional. En sus inicios, AWS tampoco tenía márgenes brutos tan altos como el software, pero al final su mercado resultó más grande que el de mucho software.

\subsection{Resumen de la sección}

El núcleo del balance de los sectores de IA en 2025 es la velocidad de los ingresos. Coding ya está validado; el video y Agentic Commerce podrían desencadenar una reorganización mayor de internet, y en las aplicaciones verticales lo determinante es el ROI de cada industria. La valuación de las aplicaciones de IA no puede basarse solo en la vieja plantilla de márgenes brutos del SaaS.
\section{¿De qué bolsillo salen los ingresos de la IA? El fondo de ingresos electrónicos y el fondo laboral}

La sección anterior analizó cada segmento; esta responde una pregunta más de fondo: ¿de qué bolsillo salen los enormes ingresos de estas empresas de IA? Freda llama a los ingresos ya existentes de internet «ingresos electrónicos/impuesto electrónico»: publicidad en línea, comisiones de comercio electrónico y suscripciones, que en conjunto rondan los 400,000 millones de dólares. Este fondo no es pequeño, pero si una sola empresa, OpenAI, quisiera llegar a 200,000 millones de dólares de ingresos, obtenerlos únicamente de este fondo implicaría que Google, Meta, Amazon y otras cedieran una parte considerable de su participación.

\begin{figure}[H]
\centering
\includegraphics[width=0.94\textwidth]{electronic-revenue-vs-labor.png}
\caption{Ingresos electrónicos frente al fondo laboral: lo que ya existe en internet no es lo bastante grande; las cifras realmente grandes están en el trabajo y la productividad. Diagrama conceptual propio, elaborado a partir de la conversación entre 01:09:51 y 01:14:11.}
\end{figure}

\begin{knowledgebox}{Lectura de la figura: por qué disputarse solo la publicidad no basta}
A la izquierda, los ingresos electrónicos incluyen publicidad, comisiones de comercio electrónico y suscripciones; el programa los sitúa en el orden de 400,000 millones de dólares. A la derecha, el fondo laboral corresponde al costo de la mano de obra en Estados Unidos, del orden de 15 billones de dólares. Si la IA solo se convierte en «un Google pequeño», su alcance es limitado; solo si entra en la atención al cliente, el trabajo del conocimiento, AI for Science y el aumento de la productividad, el margen de ingresos crecerá de verdad.
\end{knowledgebox}

\subsection{¿Crecerán los ingresos publicitarios gracias a la IA?}

Freda no está de acuerdo con la idea de que «la IA hará crecer de forma notable el gasto publicitario total». Los ingresos publicitarios de Estados Unidos han crecido de manera relativamente estable en los últimos veinte años, no hay razón para que la proporción media de los ingresos que las empresas destinan a publicidad aumente de pronto y de forma marcada, y la penetración de la publicidad en línea ya es muy alta. Por ello, que la IA se dispute la publicidad se parece más a una redistribución de lo existente que a la creación, de la nada, de un gran fondo nuevo.

\begin{warningbox}{Hay que distinguir entre crecimiento publicitario y participación publicitaria}
La IA puede cambiar quién se queda con los ingresos publicitarios, pero no necesariamente amplía de forma notable el total del mercado publicitario. La competencia entre ChatGPT, Google, TikTok, Amazon y Meta es, sobre todo, una redistribución de participación y de puntos de acceso.
\end{warningbox}

\subsection{Fondo laboral, productividad e impuesto a la IA}

Si la capacidad de los modelos mejora, las cifras realmente grandes están en el costo de la mano de obra. El PIB de Estados Unidos ronda los 30 billones de dólares y el costo de la mano de obra, unos 15 billones; solo la atención al cliente es un mercado de cientos de miles de millones de dólares. Si la IA eleva un 10\% la productividad mundial, puede generar un enorme incremento del PIB. Freda lo ilustra con un ejemplo sencillo: 100 personas producen 100 bienes; con IA, 80 personas producen 110 bienes más baratos y mejores. El PIB real crece, y las utilidades de la empresa y los ingresos de los empleados que permanecen pueden aumentar, pero las 20 personas sustituidas pasarán por una transición dolorosa.

Por eso el programa también menciona la posibilidad de un impuesto a la IA o un impuesto a los despidos. No se trata de alarmismo, sino de un problema de distribución habitual en las revoluciones de la productividad: que la producción total crezca no significa que todos se beneficien a corto plazo.

\begin{importantbox}{Las dos etapas de los ingresos de la IA}
Al principio, cuando la capacidad de los modelos todavía no ha madurado del todo, a la IA le resulta más fácil disputarse los ingresos electrónicos: publicidad, suscripciones, comisiones de comercio electrónico y presupuestos de software. Más adelante, si los Agent y AI for Science maduran, la IA podrá entrar de verdad en el fondo laboral y en el fondo de productividad adicional.
\end{importantbox}

\subsection{Resumen de la sección}

Los ingresos de la IA provienen de dos niveles. Los ingresos ya existentes de internet pueden sostener la comercialización inicial, pero no bastan para explicar inversiones del orden del billón de dólares; las cifras realmente grandes están en el costo de la mano de obra, el aumento de la productividad y la creación de nuevo valor. Para juzgar si hay una burbuja, hay que ver si la IA logra pasar del primer nivel al segundo.
\section{¿Estamos en una AI Bubble?: un tablero de indicadores, no una consigna}

La sección anterior desglosó con claridad el conjunto de ingresos; esta sección regresa a la pregunta sobre la burbuja que da título al episodio. La respuesta de Freda es: hoy no es una bubble, pero en el futuro podría convertirse en una. Rechaza discutir a partir de sentimientos y propone examinar dos niveles: primero, si hoy ya hay una burbuja; segundo, si en el futuro se inflará hasta convertirse en una. Los criterios de juicio incluyen la velocidad de adopción de las aplicaciones de IA, la materialización de los ingresos, el ROIC de las grandes empresas, si los modelos siguen mejorando y si los ingresos de IA continúan creciendo.

\begin{figure}[H]
\centering
\includegraphics[width=0.96\textwidth]{ai-bubble-dashboard.png}
\caption{Tablero de indicadores para evaluar la AI Bubble: usuarios, ingresos, ROIC, avance de los modelos y valuación. Diagrama conceptual de elaboración propia, basado en la conversación de 01:14:11--01:16:31.}
\end{figure}

\begin{knowledgebox}{Lectura de la figura: primero los indicadores verificables}
La velocidad de penetración entre usuarios indica si el producto se ha adoptado de verdad; los ingresos de OpenAI/Anthropic indican si la comercialización se ha materializado; el ROIC de las grandes empresas indica si la inversión tiene retorno; que los modelos sigan escalando determina la narrativa de largo plazo; la valuación, en cambio, debe leerse junto con estos indicadores y no puede tomarse por sí sola como evidencia de una burbuja.
\end{knowledgebox}

\subsection{Por qué considera que hoy todavía no es una bubble}

Esta subsección examina primero el momento «actual». Los indicadores de juicio ya se organizaron como un tablero; ahora la pregunta es: ¿muestran ya estos indicadores una narrativa pura que se adelanta a la realidad?

Las razones que se dan en el episodio incluyen: en menos de tres años desde la aparición de GPT, el grado de adopción entre usuarios ya equivale a lo que internet logró en diez años de desarrollo; la IA ya genera decenas de miles de millones de dólares en ingresos; tras invertir en IA, el ROIC trimestral de las grandes empresas sigue aumentando; OpenAI y Anthropic son los actores principales que pueden seguirse, y sectores verticales como Coding, generación de video, audio, atención al cliente, servicios legales y salud también ya registran ingresos. Es decir, la IA actual al menos no es pura narrativa: ya tiene usuarios, ingresos y efectos financieros.

\subsection{¿Se convertirá en una bubble en el futuro?}

A continuación, la pregunta se traslada del presente al futuro. Aunque hoy no sea una burbuja, podría convertirse en una si los ingresos, el avance de los modelos y la valuación se desacoplan entre sí.

Para saber si en el futuro habrá una burbuja, Freda observa dos variables. Primero, si los modelos siguen mejorando. Considera la aparición de Gemini 3 como una señal de que el preentrenamiento todavía puede seguir escalando, y también vale la pena observar la siguiente generación de modelos entrenada con Blackwell. Segundo, si los ingresos de IA siguen creciendo. Si los ingresos no siguen el ritmo del CAPEX, de la valuación y de las expectativas del mercado, el riesgo de burbuja aumentará.

\begin{warningbox}{Que no sea una bubble no significa que no haya empresas sobrevaluadas}
Que un gran ciclo sea real no implica que la valuación de todas las empresas sea razonable, ni que todos los segmentos vayan a tener éxito. Los ferrocarriles, internet y la computación en la nube fueron revoluciones reales, pero en cada ola también hubo una gran cantidad de empresas fallidas y un gasto de capital excesivo.
\end{warningbox}

\subsection{Resumen de la sección}

Conviene convertir el juicio sobre la burbuja en un tablero de indicadores: observar usuarios, ingresos, ROIC, avance de los modelos y valuación, en lugar de fijarse solo en el sentimiento del mercado. La IA actual ya tiene ingresos y aplicaciones reales, pero que en el futuro haya una burbuja depende de que los ingresos sigan materializándose y de que los modelos sigan mejorando.
\section{Perspectivas para 2026: de la narrativa de inversión a la materialización de ingresos}

Esta sección lleva a 2026 las conclusiones anteriores sobre los modelos, la tecnología financiera, el conjunto de ingresos y la evaluación de la burbuja. El cambio central es que el mercado empieza a exigir que los ingresos por IA se materialicen, en lugar de solo escuchar a las empresas decir que seguirán invirtiendo.

La última sección lleva a 2026 todos los marcos anteriores. Freda considera que el mercado ya cambió en los dos últimos trimestres: no basta con invertir en IA, se necesitan ingresos reales por IA. La atención del mercado pasa gradualmente de «quién compró cuántas GPU» a «quién incorpora la IA en sus ingresos, sus utilidades y su productividad».

\begin{figure}[H]
\centering
\includegraphics[width=0.96\textwidth]{ai-market-2026-watchlist.png}
\caption{Lista de seguimiento del mercado para 2026: Google, Meta, Tesla, Cloud, Nvidia/ASIC y la mejora de eficiencia con IA en las industrias tradicionales. Diagrama conceptual de elaboración propia, basado en la conversación de 01:16:31--01:24:26.}
\end{figure}

\begin{knowledgebox}{Lectura de la figura: el eje de 2026 son los ingresos y el desempeño financiero}
En Google hay que observar la competencia publicitaria y las nuevas formas de monetizar la IA; en Meta, si su modelo entra en el primer nivel; en Tesla, FSD/Robotaxi; en Cloud, la competencia entre crecimiento y margen de utilidad; en Nvidia/ASIC, si los chips de diseño propio se entregan; en las industrias tradicionales, si la mejora de eficiencia con IA aparece en los estados financieros.
\end{knowledgebox}

\subsection{Grandes empresas: Google, Meta, Tesla, Cloud, Nvidia y Apple}

La lógica de Google en el corto plazo es favorable: sus modelos están en el SOTA y tiene una buena posición estratégica, pero la competencia publicitaria se intensificará: el crecimiento del volumen publicitario de OpenAI y el margen que tiene TikTok en Estados Unidos, con mucho tiempo de uso y baja monetización, obligarán a Google a buscar nuevos modelos de negocio de IA. Meta y Tesla tienen en común que sus fundamentales no son perfectos, pero el precio de sus acciones depende en gran medida del avance de la IA. En Meta hay que observar si su modelo llega al primer nivel; en Tesla, si logra sacar del vehículo al conductor de seguridad, es decir, el avance sustancial de Robotaxi/FSD.

El crecimiento de los proveedores de nube podría ser muy bueno el próximo año, pero el panorama competitivo empeora. Los tres grandes proveedores de nube tenían márgenes muy altos; ahora se suman Oracle, Neo Cloud y las nubes propias de las empresas de modelos, y la homogeneización del alquiler de GPU y la concentración de clientes presionarán los márgenes. En Nvidia y en el sector de chips hay que observar los ASIC, los chips de diseño propio, la venta externa de las TPU de Google, el desempeño de los modelos entrenados con Blackwell y la estructura de costos del cómputo en un mundo con escasez de electricidad. Apple es la excepción: no tiene una narrativa evidente de IA y aun así su acción se mantiene fuerte; las razones podrían estar en su cultura de producto, su orientación al usuario y las expectativas sobre el relevo en la dirección.

\subsection{Posiciones largas en los beneficiarios de la IA, posiciones cortas en los perjudicados}

Esta subsección traduce las previsiones para 2026 a una estructura de posiciones largas y cortas. Como antes se dijo que el mercado entra en un periodo de materialización de ingresos, tanto los beneficiarios como los perjudicados serán más claros, y la divergencia importará más que un alza generalizada.

Freda considera que 2026 podría ser un mercado típico de fondos de cobertura, con oportunidades tanto en posiciones largas como cortas. Entre los sectores beneficiados están los semiconductores, la energía, la nube, las aplicaciones de IA y las empresas líderes de las industrias tradicionales; entre los perjudicados podrían estar la subcontratación de TI, la producción de contenido y otros eslabones intensivos en mano de obra que la IA sustituye directamente. Ella da especial importancia a los datos de mejora de eficiencia con IA de las empresas líderes de las industrias tradicionales, porque esos datos ya empiezan a aparecer en los estados financieros.

\begin{knowledgebox}{Ejemplos de mejora de eficiencia con IA en industrias tradicionales}
\begin{tabularx}{\textwidth}{p{0.22\textwidth}p{0.36\textwidth}X}
\toprule
Empresa/industria & Mejora de eficiencia mencionada en el programa & Lección \\
\midrule
Walmart & Pronóstico de demanda en la cadena de suministro, menor tasa de desabasto, menos kilómetros recorridos por los camiones & La IA entra directamente en los costos operativos y en la eficiencia de inventario. \\
Sector financiero & Aumento de montos de crédito y tasas de aprobación, menor tiempo de procesamiento, menor tasa de incumplimiento & La IA entra en las decisiones de control de riesgos y de aprobación. \\
CH Robinson & La IA se encarga de la mitad de las reservaciones; aumento de productividad & Los flujos de trabajo logísticos pueden automatizarse con IA. \\
Empresas del S\&P 500 & Por primera vez presentan datos concretos de mejora de eficiencia con IA & La IA pasa de la narrativa a los estados financieros. \\
\bottomrule
\end{tabularx}
\end{knowledgebox}

\subsection{Referencias prácticas para el público general}

Al final, el programa ofrece una referencia práctica: hoy el mercado bursátil estadounidense es un mercado de IA, y su indicador de tendencia ya no es solo Nvidia, porque lo que más le importa al mercado son los ingresos por IA, no solo la inversión en IA. El ancla más cuantificable son los ingresos de OpenAI: el crecimiento de usuarios, el crecimiento de ingresos y si el gobierno de Estados Unidos participa en sus proyectos de infraestructura influirán en la evaluación que hace el mercado del ciclo de la IA.

\begin{importantbox}{Observación central para 2026}
Si 2023--2025 fue la etapa en que «se validó la inversión en IA», 2026 se parecerá más a la etapa en que «se validan los ingresos por IA». El mercado premiará a las empresas capaces de incorporar la IA en sus ingresos, utilidades, productividad y estados financieros.
\end{importantbox}

\subsection{Resumen de la sección}

La observación de inversiones y de la industria en 2026 debe pasar del CAPEX a la materialización de ingresos. En las grandes empresas hay que observar la monetización de la IA y los márgenes; en las aplicaciones de IA, el ROI real; en las industrias tradicionales, los datos de mejora de eficiencia en los estados financieros; y en el plano macroeconómico, las elecciones intermedias, el consumo y el estímulo fiscal.
\section{Términos clave: índice de conceptos de este episodio}

Esta sección reúne los términos de finanzas, IA y mercados de capitales que aparecieron dispersos a lo largo del texto, para que el lector no se quede solo con los nombres de las empresas sin haber asimilado los conceptos. Las explicaciones de la tabla sirven al hilo conductor de este episodio y no sustituyen a textos especializados de finanzas o de aprendizaje automático.

\begin{longtable}{p{0.24\textwidth}p{0.34\textwidth}p{0.33\textwidth}}
\toprule
Término & Explicación en una frase & Papel en este episodio \\
\midrule
Altimeter Capital & Fondo tecnológico de Silicon Valley que invierte en los mercados primario y secundario & Institución de inversión del invitado y origen de su perspectiva. \\
Crossover & Estructura de fondo que invierte tanto en el mercado primario como en el secundario & Explica por qué se puede expresar una tesis sobre la IA a través de Nvidia. \\
CAPEX & Gasto de capital & Inversión de gran escala que requieren por igual los centros de datos de IA, las GPU, el entrenamiento y la reindustrialización. \\
Re-industrialization & Reindustrialización: regreso de la manufactura y la infraestructura a territorio estadounidense & Relacionada con el circuito cerrado de centros de datos, energía, aranceles y capacidad de cómputo para IA. \\
Digitization of Finance & Digitalización de la industria financiera & Contexto común de las stablecoins, los pagos, Robinhood, los mercados de predicción y el Agent Commerce. \\
Stablecoin & Moneda estable: token generalmente anclado a una moneda fiduciaria como el dólar & Posible infraestructura para los pagos de Agents y la digitalización de las finanzas. \\
Negative flywheel & Efecto bola de nieve negativo & Describe la etapa en la que los ingresos de una empresa de modelos no alcanzan a cubrir el costo de entrenar la siguiente generación. \\
Scaling Law & Regularidad empírica entre el tamaño del modelo, los datos, la capacidad de cómputo y el desempeño & Explica por qué las empresas de modelos siguen comprando GPU y entrenando. \\
API & Interfaz para invocar externamente las capacidades de un modelo & Uno de los pilares de ingresos de OpenAI. \\
Agent Commerce & Un Agent busca, compara, compra y paga en nombre del usuario & Podría redefinir los puntos de entrada de la publicidad, el comercio electrónico y los pagos. \\
Alpha & Rendimiento generado por la propia empresa, por encima de la Beta del mercado & La mejora del negocio de Robinhood se usa para ilustrar la Alpha. \\
Beta & Rendimiento que aporta el ciclo de la industria o del mercado & Las criptomonedas, la IA y el mercado tecnológico en su conjunto pueden aportar Beta. \\
Tokenización & Proceso de representar un activo como un token negociable & Una de las vías de Robinhood para entrar a Europa y para la innovación financiera. \\
Prediction Market & Mercado de predicción: usa precios de negociación para expresar la probabilidad de un evento & Se describe como una especie nueva y una muestra de la digitalización de las finanzas. \\
ROIC & Return on Invested Capital, rendimiento sobre el capital invertido & Permite juzgar si la inversión en IA de las grandes empresas tiene retorno. \\
ASIC & Application-Specific Integrated Circuit, chip de propósito específico & Dirección con la que las nubes y las empresas de modelos reducen a largo plazo el riesgo de homogeneización de las GPU. \\
Neo Cloud & Proveedores de servicios en la nube de nueva generación para GPU/IA & Ejemplo de cómo la competencia en la nube se amplía de las tres grandes nubes a más participantes. \\
AI Tax & Propuesta de redistribución para atender los problemas distributivos causados por la sustitución de trabajo por IA & Corresponde al problema social entre el aumento de la productividad y el impacto en el empleo. \\
\bottomrule
\end{longtable}

\subsection{Resumen de la sección}

Los términos de este episodio abarcan finanzas, IA, internet y mercados de capitales. Entender la relación entre ellos permite convertir el programa de «rumores del mundo de la inversión» en una clase sobre estructura industrial: cómo interactúan el gasto de capital, la bolsa de ingresos, el control de los puntos de entrada, la ejecución organizacional y la productividad.
\section{Síntesis y ampliación}

\subsection{Conclusiones centrales}

Esta sección presenta primero las conclusiones de todo el texto y después deja planteadas las preguntas abiertas. Así, el lector puede condensar este episodio, que pasa de ser una «entrevista sobre mercados de capitales» a un marco de trabajo para evaluar el ciclo de la IA.

Por último, el episodio se condensa en diez juicios.

\begin{enumerate}
\item La línea principal de la tecnología en la bolsa estadounidense en 2025--2026 es la confluencia de tres corrientes: la IA, la reindustrialización y la digitalización financiera.
\item La tesis clave de OpenAI que va contra el consenso es que se le ve como una empresa de productos, y no solo como una empresa de modelos.
\item En el corto plazo, el flujo de efectivo de las empresas de modelos se comporta como una bola de nieve negativa; el punto de inflexión llega cuando baja el ritmo de crecimiento de los costos de entrenamiento o cuando sube el ROI de los modelos.
\item Los ingresos de OpenAI no deben evaluarse solo por las suscripciones individuales; también hay que considerar el segmento empresarial, la API, los Agent, la publicidad y el comercio electrónico.
\item La divergencia entre Anthropic y OpenAI se debe a la estructura de clientes, la ruta de producto y los supuestos de sus proyecciones financieras.
\item Robinhood muestra cómo un negocio cíclico puede generar Alpha mediante la diversificación, la participación de mercado, el poder de fijación de precios y el control de costos.
\item El «chico malo» que prefiere el capital de Silicon Valley no es caprichoso: encarna una capacidad de ejecución founder-led que va contra la intuición.
\item Coding, video, Agentic Commerce y las aplicaciones verticales son los segmentos de aplicaciones de IA de 2025 que más vale la pena estudiar.
\item En su etapa inicial, los ingresos de la IA le disputan su parte a la bolsa de ingresos de la electrónica; a largo plazo, deben entrar en la bolsa del trabajo y en la nueva productividad.
\item Para juzgar si hay una AI bubble hay que observar usuarios, ingresos, ROIC, avance de los modelos y valuación, no solo el sentimiento del mercado.
\end{enumerate}

\subsection{Preguntas abiertas}

Al final se dejan planteadas preguntas abiertas porque muchas de las variables discutidas en este episodio siguen evolucionando: los ingresos de los modelos, el Agent Commerce, los puntos de entrada financieros, los márgenes de la nube y el reflejo financiero de las ganancias de eficiencia de la IA requieren verificarse en los próximos trimestres.

\begin{itemize}
\item ¿El punto de entrada empresarial de OpenAI rendirá frutos antes que el superpunto de entrada para consumidores?
\item ¿Podrá la ruta 2B/Coding de Anthropic mantener su ritmo de crecimiento y lograr una economía unitaria mejor que la de OpenAI?
\item ¿El Agentic Commerce perjudicará primero a las grandes plataformas o beneficiará primero a los comerciantes de cola larga?
\item ¿Podrá la aplicación financiera integral de Robinhood entrar en el mercado masivo de la gestión patrimonial y los activos alternativos?
\item ¿El criterio de valuación de las empresas de aplicaciones de IA volverá al margen bruto o se desplazará hacia el margen bruto por usuario y la utilidad absoluta?
\item ¿Se convertirán en 2026 los ingresos de la IA en un ancla de mercado más importante que los envíos de Nvidia?
\end{itemize}

\subsection{Lecturas adicionales}

\begin{itemize}
\item EP127, informe trimestral de modelos grandes de fin de año de Guang Mi (广密): AI War, Online Learning y la pila de ingresos.
\item EP136, episodio 9 del informe trimestral de modelos grandes de Guang Mi: Coding, el OS del modelo y los tres grandes de Silicon Valley.
\item EP139, panorama de los Agent: historia de la tecnología de Agent, el OpenClaw Moment y su impacto social.
\item Entrevistas públicas relacionadas de Altimeter/BG2: Sam Altman, Jensen Huang y el relato del capital tecnológico estadounidense.
\item Materiales públicos de OpenAI, Anthropic, Robinhood, Waymo, Polymarket/Kalshi, entre otros: sirven para contrastar los cambios en ingresos, productos y regulación.
\end{itemize}

\begin{importantbox}{El juicio final}
Lo que más vale la pena conservar de EP125 no es una cifra de valuación concreta, sino una forma de leer los mercados de capitales: primero preguntar por la bolsa de ingresos y después por la curva de costos; primero distinguir entre Beta y Alpha y después evaluar la ejecución de la empresa; primero ver si la inversión en IA genera ingresos y ROIC, y solo entonces discutir la burbuja. Leído así, el ciclo de la IA es menos propenso a dejarse arrastrar por un relato único.
\end{importantbox}

\end{document}
